\textbf{Phần 1. Trắc nghiệm lựa chọn}
\setcounterpageref{ex}{0}

\begin{ex} %[3D2V42]
Tính giá trị của biểu thức: $M = 350 - 50 \times 4$.
\choice
{$1200$}
{$200$}
{$250$}
{\True $150$}
\loigiai{
Nhân trước: $50 \times 4 = 200$. Trừ sau: $350 - 200 = 150$.
\textbf{Đáp án của câu hỏi là D}
}
\end{ex}

\begin{ex} %[3D2V42]
Tính giá trị của biểu thức có dấu ngoặc: $N = (350 - 50) : 5$.
\choice
{$340$}
{\True $60$}
{$50$}
{$70$}
\loigiai{
Trong ngoặc: $350 - 50 = 300$. Chia sau: $300 : 5 = 60$.
\textbf{Đáp án của câu hỏi là B}
}
\end{ex}

\begin{ex} %[3D2V42]
Thực hiện phép tính: $48 : 4 \times 2$.
\choice
{$6$}
{$12$}
{\True $24$}
{$16$}
\loigiai{
Chỉ có nhân chia, làm từ trái sang phải: $48 : 4 = 12$; $12 \times 2 = 24$.
\textbf{Đáp án của câu hỏi là C}
}
\end{ex}

\begin{ex} %[3D2V42]
Biểu thức nào có giá trị bằng 100?
\choice
{\True $400 : 4 + 25 - 25$}
{$400 : (4 + 25) - 25$}
{$400 - 4 \times 25$}
{$400 : 2 - 50$}
\loigiai{
$400 : 4 + 25 - 25 = 100 + 0 = 100$.
\textbf{Đáp án của câu hỏi là A}
}
\end{ex}

\begin{ex} %[3D2V42]
Tìm $x$, biết: $240 : (x - 5) = 8$.
\choice
{$30$}
{$40$}
{$25$}
{\True $35$}
\loigiai{
$x - 5 = 240 : 8 = 30 \Rightarrow x = 30 + 5 = 35$.
\textbf{Đáp án của câu hỏi là D}
}
\end{ex}

\begin{ex} %[3D2V42]
Tính giá trị biểu thức: $125 \times 4 - 300 : 3$.
\choice
{$350$}
{\True $400$}
{$450$}
{$500$}
\loigiai{
$125 \times 4 = 500$; $300 : 3 = 100$; $500 - 100 = 400$.
\textbf{Đáp án của câu hỏi là B}
}
\end{ex}

\begin{ex} %[3D2V42]
Đặt dấu ngoặc thích hợp để biểu thức $15 + 35 : 5$ có giá trị bằng 10.
\choice
{$15 + (35 : 5)$}
{$(15 + 35 : 5)$}
{\True $(15 + 35) : 5$}
{Không thể đặt được}
\loigiai{
$(15 + 35) : 5 = 50 : 5 = 10$.
\textbf{Đáp án của câu hỏi là C}
}
\end{ex}

\begin{ex} %[3D2V42]
Tính nhanh giá trị biểu thức: $45 \times 6 + 45 \times 4$.
\choice
{\True $450$}
{$400$}
{$500$}
{$480$}
\loigiai{
$45 \times (6 + 4) = 45 \times 10 = 450$.
\textbf{Đáp án của câu hỏi là A}
}
\end{ex}

\begin{ex} %[3D2V42]
Một người có 3 giỏ cam, mỗi giỏ có 15 quả cam và 2 giỏ táo, mỗi giỏ có 20 quả táo. Biểu thức tính tổng số quả cả hai loại là
\choice
{$(15 + 20) \times (3 + 2)$}
{$15 + 3 \times 20 + 2$}
{$15 \times 20 + 3 \times 2$}
{\True $15 \times 3 + 20 \times 2$}
\loigiai{
Số cam là $15 \times 3$, số táo là $20 \times 2$. Tổng số quả: $15 \times 3 + 20 \times 2$.
\textbf{Đáp án của câu hỏi là D}
}
\end{ex}

\begin{ex} %[3D2V42]
Tính giá trị biểu thức: $P = 500 - 100 : (2 + 3)$.
\choice
{$80$}
{\True $480$}
{$450$}
{$490$}
\loigiai{
Trong ngoặc: $2 + 3 = 5$. Chia: $100 : 5 = 20$. Trừ: $500 - 20 = 480$.
\textbf{Đáp án của câu hỏi là B}
}
\end{ex}

\begin{ex} %[3D2V42]
Tìm số tự nhiên $y$ sao cho: $(y + 15) \times 3 = 90$.
\choice
{$10$}
{$20$}
{\True $15$}
{$25$}
\loigiai{
$y + 15 = 90 : 3 = 30 \Rightarrow y = 30 - 15 = 15$.
\textbf{Đáp án của câu hỏi là C}
}
\end{ex}

\begin{ex} %[3D2V42]
So sánh hai biểu thức: $A = (40 + 20) \times 2$ và $B = 40 + 20 \times 2$.
\choice
{\True $A > B$}
{$A < B$}
{$A = B$}
{Không so sánh được}
\loigiai{
$A = 60 \times 2 = 120$. $B = 40 + 40 = 80$. Vì $120 > 80$ nên $A > B$.
\textbf{Đáp án của câu hỏi là A}
}
\end{ex}

\textbf{Phần 2. Trắc nghiệm đúng sai}
\setcounterpageref{ex}{0}

\begin{ex} %[3D2V42]
Xét tính đúng sai của việc tính giá trị biểu thức:
\choiceTF
{\True Biểu thức $100 - 40 + 10$ có giá trị bằng $70$.}
{Biểu thức $100 - (40 + 10)$ có giá trị bằng $70$.}
{\True Dấu ngoặc thay đổi thứ tự ưu tiên của các phép tính trong biểu thức.}
{Tất cả các biểu thức chứa phép nhân đều phải làm từ phải sang trái.}
\loigiai{
\begin{itemize}
    \item Ý a) Đúng: $100 - 40 = 60$; $60 + 10 = 70$.
    \item Ý b) Sai: Sai vì $100 - 50 = 50$.
    \item Ý c) Đúng: Bản chất toán học của ngoặc đơn.
    \item Ý d) Sai: Sai vì nhân chia làm từ trái sang phải.
\end{itemize}
\begin{itemize}
    \item a) \textbf{Đúng}.
    \item b) \textbf{Sai}.
    \item c) \textbf{Đúng}.
    \item d) \textbf{Sai}.
\end{itemize}
}
\end{ex}

\begin{ex} %[3D2V42]
Cho biểu thức $T = 72 : 8 \times 3$:
\choiceTF
{\True Thực hiện phép tính $72 : 8$ trước.}
{\True Kết quả của $72 : 8$ là 9.}
{\True Giá trị của cả biểu thức $T$ bằng 27.}
{Nếu làm phép nhân $8 \times 3$ trước thì giá trị biểu thức không thay đổi.}
\loigiai{
\begin{itemize}
    \item Ý a) Đúng: Quy tắc từ trái sang phải.
    \item Ý b) Đúng: $72 : 8 = 9$.
    \item Ý c) Đúng: $9 \times 3 = 27$.
    \item Ý d) Sai: Sai vì $72 : 24 = 3 \neq 27$.
\end{itemize}
\begin{itemize}
    \item a) \textbf{Đúng}.
    \item b) \textbf{Đúng}.
    \item c) \textbf{Đúng}.
    \item d) \textbf{Sai}.
\end{itemize}
}
\end{ex}

\textbf{Phần 3. Trả lời ngắn}
\setcounterpageref{ex}{0}

\begin{ex} %[3D2V42]
Tính giá trị của biểu thức: $180 - 80 : 4 \times 2$.
\shortans{140}
\loigiai{
\textbf{Đáp số:} $140$\\
$80 : 4 = 20$; $20 \times 2 = 40$; $180 - 40 = 140$.
}
\end{ex}

\begin{ex} %[3D2V42]
Tính giá trị của biểu thức: $(150 + 250) : (4 + 4)$.
\shortans{50}
\loigiai{
\textbf{Đáp số:} $50$\\
$400 : 8 = 50$.
}
\end{ex}

\begin{ex} %[3D2V42]
Tìm $x$, biết: $x + 25 \times 4 = 250$.
\shortans{150}
\loigiai{
\textbf{Đáp số:} $150$\\
$x + 100 = 250 \Rightarrow x = 150$.
}
\end{ex}

\begin{ex} %[3D2V42]
Tính nhanh giá trị biểu thức: $38 \times 5 + 38 \times 5$.
\shortans{380}
\loigiai{
\textbf{Đáp số:} $380$\\
$38 \times 10 = 380$.
}
\end{ex}

\begin{ex} %[3D2V42]
Một người mua 5 gói kẹo giá 12 nghìn đồng mỗi gói và đưa cho cô bán hàng tờ 100 nghìn đồng. Hỏi người đó nhận lại bao nhiêu nghìn đồng tiền thừa?
\shortans{40}
\loigiai{
\textbf{Đáp số:} $40$\\
$100 - 12 \times 5 = 100 - 60 = 40$ nghìn đồng.
}
\end{ex}

\begin{ex} %[3D2V42]
Tìm số tự nhiên $k$ biết: $60 : k + 15 = 25$.
\shortans{6}
\loigiai{
\textbf{Đáp số:} $6$\\
$60 : k = 10 \Rightarrow k = 6$.
}
\end{ex}
